\documentclass{article}
\usepackage[hidelinks]{hyperref}
\title{A supplier change before a native station effect has settled}
\begin{document}
\maketitle

\paragraph{Envelope and law.}
This follow-up to
\href{https://github.com/eternal-november/eos-research-programmes/pull/634}{the durable release replay}
moves the second fault before the actual station effect in three admitted
crash-phase branches. The duty is exactly one named native station effect
by cut 7; the source must stay off its bay through that effect and return
by cut 7. Independent warehouse motions and radio callbacks are due at
cuts 3 and 6. The protected carrier can return in one native step.
An intended primary station command arrives through native radio at
cut 2. Its worker then exits before publication or after publication
without a reply. The new primary request is tried at cut 7. A second
fault can make P's path unusable before that retry. A qualified Q can
retry the unchanged operation. The receiver has a single reachable
SQLite authority for the native effect; neither router sees the
evaluator's hidden path or old paid flag before it commits to its
preparation decision. Synchronized outer cuts, fail-stop phases,
echoed native callbacks, setup semantics and trusted authorizations are
fixture premises. This is not an independent mission.

The positive envelope requires either the primary retry path to survive,
or both a real setup delivery to Q at cut 4 and Q's retry path to work
at cut 7. A reply must reach the source at cut 7. Both normal and
compound branches are admitted; the primary can be lost on histories
where no station effect yet exists. Outside the stated branches, refusal
to reenter is safe but does not count as successful service. Common
cause loss of both paths and loss of the confirmation are exercised as
boundary cases, not counted in the guaranteed service envelope.

\paragraph{Concrete interfaces and proof.}
Setup sends two bytes on native radio; a callback equips Q with named
retry authority after two fixture ticks. P and Q each submit the same
three-byte operation, and the receiver checks the exact app key against
the pinned local atomic step journal. A native erased send still spends
sender bytes. The old cut-2 debit has expired before cut 7, so healthy P
can retry. If P loses the path at cut 7, its debit is on P only;
the broker gives no accepted request and did not charge receiver capacity.
The preauthorized Q then fits into both its own ledger and the common
receiver ledger, and invokes the same journal key. Side work uses its
own accepted destination and source motion. Its one-byte transmissions
and the setup can share the Side transmitter while respecting the
five-byte limit; the return receipt and side traffic fit at their
common destination. Neither a prepared sender nor a transmitted request
is itself a station effect.

Each delivered eligible request either creates the durable native
transition or returns its stored result. Induction over delivered
attempts preserves at most one published effect, including a late
primary request after the source returns. At least one such attempt
arrives by cut 7 in every admitted positive branch by the route and
capacity argument. The qualified response is emitted only after that
result becomes authoritative. The return follows its actual callback;
if the callback is erased, the carrier remains held. Hence the joint
guarantee follows through the same actual histories in which both
independent jobs run, without treating their success as evidence of the
paid key.

\paragraph{Comparison and scope.}
A reactive policy waits to see P's actual failure, then starts the
same setup at cut 7. It cannot finish the two-tick preparation by the
required cut; it misses even if Q is healthy. The lower-bound witness
fixes the smallest setup step budget and two histories with identical
preparation views and opposite later primary path. Preparing only
after observing P's later loss cannot succeed on the lost-P case.
An equally informed direct RIFL-style application with pre-staged
backup and the same route effects ties; the reactive arm is only
a timing ablation, not the strongest established composite system.
Pre-staging costs a real setup message even when the primary is healthy.

The held-in sweep has 64 configurations per arm. The staged and direct
policies each meet all 20 cases in the positive envelope; reactive meets
16. In the other cases all three preserve exactly-once station effect
and safe source return by the tested late cut. The counted request,
setup and reply bytes total 1,004 for staging versus 880 for reactive;
that comparison includes late attempts and excludes the identical
independent traffic. These are finite schedule totals, not estimates.

Every scenario also tries a returning P at cut 12. Its actual callback
cannot duplicate the effect. A missing service case can settle and
return safely at that later cut in the fixture, missing the bound.
The replay runs the native station worker and both native plant/radio
adapters end to end for each arm. These checks establish a reusable
held-in envelope and explicit local readiness conditions, not inventory
coverage, field validity, a frozen transfer, or advancement over the
strongest published assembly.
\end{document}
